\documentclass{article}
\usepackage{graphicx} % Required for inserting images
\usepackage[table]{xcolor}
\definecolor{azulCeleste}{HTML}{00ADEF}

\title{refetencias}
\author{cristopher diaz}
\date{May 2026}

\begin{document}
    
\maketitle

\section{Referencias}
\centering
    \renewcommand{\arraystretch}{1.8} % Le da altura y elegancia a las celdas
    \setlength{\tabcolsep}{12pt}      % Espacio entre columnas
    
    \begin{tabular}{lllll} % 'l' significa alineado a la izquierda (left)
        \rowcolor{azulCeleste} % Pinta la fila de azul
        \textbf{\color{white}Instrucción} & 
        \textbf{\color{white}Ejemplo} & 
        \textbf{\color{white}Significado} & 
        \textbf{\color{white}Comentario} \\
        
        % Aquí empieza el contenido (puedes borrar estas líneas si solo quieres la barra)
        add & add \$s1,\$s2,\$s3  & \$sl = \$s1 = \$s2  & tres operando: dato en registro \\
        addi & addi \$s1,\$s2,\$s3  & \$s1 = \$s2 + 100 &  usado para sumar constantes \\
        lw & lw \$s1,100(\$s2)  & \$s1 = Memory[\$s2 + 100] &  memoria de registro \\
        sw & sw \$s1,100(\$s2)  &  Memory[\$s2 + 100] =\$s1 &  media palabra de rgistro de memoria\\
        sll & sll \$s1,\$s2,10  & \$s1 = \$s2 << 10 &  desplazamiento de la constante a la izquierda \\
         srl & srl \$s1,\$s2,10  & \$s1 = \$s2 >> 10 &  desplazamiento de la constante a la derecha\\
        beq & beq \$s1,\$s2,L  & if(\$s1 == \$s2)go to L PC+4+100 &  comprueba la bifurcacion e igualdad relativo al pc \\
        bne & bne \$s1,\$s2,L  & if(\$s1 != \$s2)go to L PC+4+100 &  se comprueba la desigualdad y bifurcacion relativa al pc \\
        slt & slt \$s1,\$s2,\$s3  & if(\$s2 < \$s3) \$s1 = 1; else \$s1=0 &  comprueba si es menor que ; usado con beq, bne \\
        slti & slti \$s1,\$s2,100  & if(\$s2 < 100 )\$s1 = 1; else \$s1 = 0 &  comparar si es menor que una constante \\
        j & j 2500  & go to 100000 &  salto de direccion destino  \\
    \end{tabular}
\end{document}
